\section{静电场中的电介质}
电介质的定义：电介质是电阻率很大，导电性能很差的物质。

电介质的极化：当电介质处在电场中时，电介质中，无论是原子中的电子还是分子中的离子或是晶体晶格上的带电粒子，在电场的作用下都会在原子大小的范围内移动。

当达到静电平衡时，在电介质表面层或体内会出现极化电荷，这个现象称作电介质的极化。

\subsection{电介质的电结构与极化}
\subsubsection{无极分子}
分子内正、负电荷中心重合，等效电偶极矩等于零，电介质整体呈电中性。

无极分子的极化来自正、负电荷中心的相对位移，所以常叫做位移极化。

\subsubsection{极性分子}
分子内正、负电荷中心不相重合。设极性分子的正负电荷中心之间的距离为$l$，分子中全部正电荷或负电荷的总电荷量为$q$，则等效电偶极矩$\vec{p}=q\vec{l}$，电介质整体呈电中性。

极性分子的极化就是等效电偶极子转向外电场的方向，叫做取向极化（一般来说，分子在取向极化的同时还会产生位移极化，但极性分子电介质以取向极化为主）。

\subsubsection{极化后电介质中的电荷分布}
如果电介质是均匀的，则在它内部仍然保持电中性，但是在电介质的两个和外电场强度$\vec{E_0}$相垂直的表面层里（厚度为分子等效电偶极矩的轴长$l$），将会分别出现正电荷和负电荷。

\subsection{电极化强度}
两类电介质极化的微观过程虽然不同，但宏观的效果却是相同的。都是在电介质的两个相对表面上出现了异号的极化电荷，在电介质内部有沿电场方向的电偶极矩。

我们取单位体积内分子电偶极矩的矢量和，即
$$\vec{P}=\frac{\sum \vec{p}}{\Delta V}$$
作为量度电介质极化程度的基本物理量，称为该点的电极化强度或$\vec{P}$矢量。单位是$C/m^2$。

\subsection{电极化强度与极化电荷的关系}
对于均匀电介质，其极化电荷只集中在表面层里或在两种不同的界面层里。

设$\vec{e_n}$为介质表面的单位法向矢量，那么
$$\sigma'=\vec{P}\cdot\vec{e_n}=P_n$$
即介质极化所产生的极化电荷面密度等于电极化强度沿介质表面外法线的分量。

\subsection{介质中的静电场}
\subsubsection{物理量}
电介质的电极化率：$\chi_e$，和电介质的性质有关，量纲为1。

电介质的相对电容率：$\varepsilon_r$

电介质的电容率/介电常量：$\varepsilon$

\subsubsection{公式}
对于各向同性电介质，有：

$\vec{E}=\vec{E_0}+\vec{E'}$

$\vec{P}=\chi_e\varepsilon_0\vec{E}$

$\varepsilon_r=1+\chi_e$

$\varepsilon=\varepsilon_r\varepsilon_0=(1+\chi_e)\varepsilon_0$

\subsubsection{（特例：平行板电容器）}
除满足以上公式外，还有

$E_0=\sigma_0/\varepsilon_0$

$E'=\sigma'/\varepsilon_0$

$E=\frac{E_0}{1+\chi_e}$

\subsubsection{拓展：非各向同性电介质}
自然界也存在一些电介质，在一定的温度范围内电容率随电场强度而变化，它们的极化规律有着复杂的线性关系。

如铁电体、永驻体、压电效应、电致伸缩等。